\documentclass[tikz,border=8pt]{standalone}
\usepackage{tikz}
\usepackage{amsmath}
\usepackage{amssymb}
\usepackage{pgfplots}
\pgfplotsset{compat=1.18}
\usepackage{ctex}
\usetikzlibrary{calc,positioning,arrows.meta,matrix}

% LC 206 反转链表
% 5 节点链表 1→2→3→4→5，展示 prev/curr/next 三指针演化

\begin{document}
\begin{tikzpicture}[
  every node/.style={font=\small},
  listnode/.style={draw=gray!70, rectangle, rounded corners=3pt, minimum width=0.85cm, minimum height=0.65cm, thick, fill=blue!8},
  nullnode/.style={font=\scriptsize, gray},
  ptr/.style={->,>=Stealth,thick},
  snap/.style={draw=gray!40, fill=white, rounded corners=4pt, inner sep=6pt, thin},
]

%% ── 标题 ──────────────────────────────────────────────────
\node[font=\large\bfseries] at (5.5, 1.2)
  {反转链表（LC 206 Reverse Linked List）};
\node[font=\small,gray] at (5.5, 0.6)
  {三指针迭代法：prev / curr / next，逐步翻转箭头方向};

%% ─────────────────────────────────────────────────────────
%% 快照 0：初始状态
%% ─────────────────────────────────────────────────────────
\def\snapay{0.0}% top y of snapshot 0

\node[font=\small\bfseries,blue!70] at (0.5, \snapay-0.25) {Step 0 初始};

% 节点
\node[listnode] (s0n1) at (1.6, \snapay-0.9) {1};
\node[listnode] (s0n2) at (3.2, \snapay-0.9) {2};
\node[listnode] (s0n3) at (4.8, \snapay-0.9) {3};
\node[listnode] (s0n4) at (6.4, \snapay-0.9) {4};
\node[listnode] (s0n5) at (8.0, \snapay-0.9) {5};
\node[nullnode] (s0null) at (9.4, \snapay-0.9) {null};

% 正向箭头
\draw[ptr,gray!60] (s0n1.east) -- (s0n2.west);
\draw[ptr,gray!60] (s0n2.east) -- (s0n3.west);
\draw[ptr,gray!60] (s0n3.east) -- (s0n4.west);
\draw[ptr,gray!60] (s0n4.east) -- (s0n5.west);
\draw[ptr,gray!60] (s0n5.east) -- (s0null.west);

% prev/curr 指针标签
\node[font=\scriptsize,green!60!black] at (0.6, \snapay-0.9) {\textbf{prev}};
\draw[->,>=Stealth,green!60!black,thin] (0.9,\snapay-0.9) -- (1.2,\snapay-0.9);
\node[font=\scriptsize,green!60!black,anchor=north] at (0.6,\snapay-0.78) {(null)};

\node[font=\scriptsize,red!70!black] at (1.6, \snapay-0.3) {\textbf{curr}};
\draw[->,>=Stealth,red!70!black,thin] (1.6,\snapay-0.46) -- (1.6,\snapay-0.62);

%% ─────────────────────────────────────────────────────────
%% 快照 1：翻转 1→2，curr 指向 2
%% ─────────────────────────────────────────────────────────
\def\snapby{-2.0}

\node[font=\small\bfseries,blue!70] at (0.5, \snapby-0.25) {Step 1};

\node[listnode] (s1n1) at (1.6, \snapby-0.9) {1};
\node[listnode] (s1n2) at (3.2, \snapby-0.9) {2};
\node[listnode] (s1n3) at (4.8, \snapby-0.9) {3};
\node[listnode] (s1n4) at (6.4, \snapby-0.9) {4};
\node[listnode] (s1n5) at (8.0, \snapby-0.9) {5};
\node[nullnode] (s1null) at (9.4, \snapby-0.9) {null};
\node[nullnode] (s1prevnull) at (0.2, \snapby-0.9) {};

% 已翻转：1←null（1 指向 null=prev）
\draw[ptr,orange!70] (s1n1.west) -- (s1prevnull.east)
  node[midway,above,font=\tiny,orange!80!black] {翻转};
% 未翻转
\draw[ptr,gray!60] (s1n2.east) -- (s1n3.west);
\draw[ptr,gray!60] (s1n3.east) -- (s1n4.west);
\draw[ptr,gray!60] (s1n4.east) -- (s1n5.west);
\draw[ptr,gray!60] (s1n5.east) -- (s1null.west);

% 指针标签
\node[font=\scriptsize,green!60!black,anchor=south] at (1.6,\snapby-1.55) {\textbf{prev}};
\draw[->,>=Stealth,green!60!black,thin] (1.6,\snapby-1.42) -- (1.6,\snapby-1.22);

\node[font=\scriptsize,red!70!black] at (3.2, \snapby-0.3) {\textbf{curr}};
\draw[->,>=Stealth,red!70!black,thin] (3.2,\snapby-0.46) -- (3.2,\snapby-0.62);

%% ─────────────────────────────────────────────────────────
%% 快照 2：翻转 2→1，curr 指向 3
%% ─────────────────────────────────────────────────────────
\def\snapcy{-4.0}

\node[font=\small\bfseries,blue!70] at (0.5, \snapcy-0.25) {Step 2};

\node[listnode] (s2n1) at (1.6, \snapcy-0.9) {1};
\node[listnode] (s2n2) at (3.2, \snapcy-0.9) {2};
\node[listnode] (s2n3) at (4.8, \snapcy-0.9) {3};
\node[listnode] (s2n4) at (6.4, \snapcy-0.9) {4};
\node[listnode] (s2n5) at (8.0, \snapcy-0.9) {5};
\node[nullnode] (s2null) at (9.4, \snapcy-0.9) {null};

% 已翻转
\draw[ptr,orange!70] (s2n2.west) to[bend left=30] (s2n1.north west);
\draw[ptr,orange!70,dashed] (s2n1.west) -- (0.2,\snapcy-0.9);
% 未翻转
\draw[ptr,gray!60] (s2n3.east) -- (s2n4.west);
\draw[ptr,gray!60] (s2n4.east) -- (s2n5.west);
\draw[ptr,gray!60] (s2n5.east) -- (s2null.west);

\node[font=\scriptsize,green!60!black,anchor=south] at (3.2,\snapcy-1.55) {\textbf{prev}};
\draw[->,>=Stealth,green!60!black,thin] (3.2,\snapcy-1.42) -- (3.2,\snapcy-1.22);

\node[font=\scriptsize,red!70!black] at (4.8, \snapcy-0.3) {\textbf{curr}};
\draw[->,>=Stealth,red!70!black,thin] (4.8,\snapcy-0.46) -- (4.8,\snapcy-0.62);

%% ─────────────────────────────────────────────────────────
%% 快照 3：最终结果 5→4→3→2→1→null
%% ─────────────────────────────────────────────────────────
\def\snapdy{-6.0}

\node[font=\small\bfseries,blue!70] at (0.5, \snapdy-0.25) {最终结果};

\node[listnode,fill=green!10,draw=green!60!black] (s3n5) at (1.6, \snapdy-0.9) {5};
\node[listnode,fill=green!10,draw=green!60!black] (s3n4) at (3.2, \snapdy-0.9) {4};
\node[listnode,fill=green!10,draw=green!60!black] (s3n3) at (4.8, \snapdy-0.9) {3};
\node[listnode,fill=green!10,draw=green!60!black] (s3n2) at (6.4, \snapdy-0.9) {2};
\node[listnode,fill=green!10,draw=green!60!black] (s3n1) at (8.0, \snapdy-0.9) {1};
\node[nullnode] (s3null) at (9.4, \snapdy-0.9) {null};

\draw[ptr,green!60!black,very thick] (s3n5.east) -- (s3n4.west);
\draw[ptr,green!60!black,very thick] (s3n4.east) -- (s3n3.west);
\draw[ptr,green!60!black,very thick] (s3n3.east) -- (s3n2.west);
\draw[ptr,green!60!black,very thick] (s3n2.east) -- (s3n1.west);
\draw[ptr,green!60!black,very thick] (s3n1.east) -- (s3null.west);

\node[font=\scriptsize,red!70!black] at (1.6, \snapdy-0.3) {\textbf{head}};
\draw[->,>=Stealth,red!70!black,thin] (1.6,\snapdy-0.46) -- (1.6,\snapdy-0.62);

%% ── 分隔线 ──────────────────────────────────────────────
\draw[gray!25, dashed] (0.0, -8.5) -- (11.0, -8.5);

%% ── 算法说明 ─────────────────────────────────────────────
\node[draw=gray!50, fill=white, rounded corners=3pt, font=\small,
      inner sep=8pt, align=left]
  at (5.5, -9.6)
  {\textbf{三指针迭代（O(n) 时间，O(1) 空间）：}\\[4pt]
   1. 初始化 prev = null，curr = head\\[2pt]
   2. 每步：next = curr.next \quad（保存后继）\\
   \quad\quad\quad curr.next = prev \quad（翻转箭头）\\
   \quad\quad\quad prev = curr \quad（prev 前移）\\
   \quad\quad\quad curr = next \quad（curr 前移）\\[2pt]
   3. curr == null 时，prev 即为新头节点};

\end{tikzpicture}
\end{document}
